\section{Исследовательская часть}
\subsection{Цель и методика исследования}

Целью исследования является анализ влияния параметров алгоритма трассировки путей
на производительность.

Исследованы две ключевые характеристики:
\begin{itemize}
    \item количество лучей на пиксель при изменении от 10 до 100 с шагом 10;
    \item максимальная глубина рекурсии при изменени от 1 до 15 с шагом 2;
\end{itemize}

Для оценки производительности измеряется время рендеринга одного кадра фиксированного разрешения~$400 \times 400$ пикселей.
Каждое измерение повторяется~100 раза, результаты усредняются.

\subsection{Технические характеристики}

Исследование проводилось на компьютере со следующими характеристиками~\cite{macbook-m3}:

\begin{itemize}
    \item \textbf{Процессор}: Apple M3, 8 ядер (4 производительных + 4 энергоэффективных);
    \item \textbf{Оперативная память}: 16~ГБ LPDDR5;
    \item \textbf{Операционная система}: macOS Tahoe 26.1 (25B78);
    \item \textbf{Компилятор}: Apple Clang 16.0.0;
\end{itemize}

Во время замеров времени ноутбук был подключён к сети электропитания.
Никакие пользовательские задачи в фоне не выполнялись. Оптимизация компилятора выключена флагом -o0~\cite{compilator-level-optimize}.

\subsection{Результаты исследований}

\subsubsection{Зависимость времени рендеринга от количества лучей}

Результаты исследования зависимости времени выполнения реализации алгоритма трассировки пути для генерации одного кадра от количества испускаемых лучей на один пиксель представлены на рисунке~\ref{fig:exp-samples}.
Глубина рекурсии была зафиксирована на десяти переотражениях.
\begin{figure}[H]
	\centering
	\includegraphics[width=1\textwidth]{img/sample_chart.png}
	\caption{график зависимости временя выполнения реализации алгоритма от количества лучей на пиксель}
	\label{fig:exp-samples}
\end{figure}


\subsubsection{Зависимость времени генерации кадра от количества переотражений}

Результаты исследования зависимости времени выполнения реализации алгоритма трассировки пути для генерации одного кадра от максимального количества переотражений одного луча представлены на рисунке~\ref{fig:exp-depth}.
Количество испускаемых лучей на один пиксель было ограничено десятью единицами.
\begin{figure}[H]
	\centering
	\includegraphics[width=1\textwidth]{img/depth_chart.png}
	\caption{график зависимости временя выполнения реализации алгоритма от максимального количества переотражений луча}
	\label{fig:exp-depth}
\end{figure}

\subsection{Вывод}

Было проведено исследование зависимости времени выполнения реализации алгоритма от его параметров, в ходе которого было выявлено:
\begin{enumerate}
    \item время генерации одного кадра линейно зависит от количества лучей на пиксель;
    \item время генерации одного кадра возрастает при увелечении глубины рекурсии: с увеличением числа допустимых переотражений наблюдается умеренный рост вычислительных затрат, однако данный эффект остаётся некритичным для общей производительности рендеринга.
\end{enumerate}

\clearpage
